\documentclass[hyperref={colorlinks=true}]{beamer} 
\usecolortheme{lily}
% we need these because of source-highlight output
\DefineNamedColor{named}{Blue}          {cmyk}{1,1,0,0}
\DefineNamedColor{named}{BrickRed}      {cmyk}{0,0.89,0.94,0.28}
\DefineNamedColor{named}{Brown}         {cmyk}{0,0.81,1,0.60}
\DefineNamedColor{named}{ForestGreen}   {cmyk}{0.91,0,0.88,0.12}
\DefineNamedColor{named}{Purple}        {cmyk}{0.45,0.86,0,0}
\DefineNamedColor{named}{Red}           {cmyk}{0,1,1,0}
\DefineNamedColor{named}{Black}         {cmyk}{0,0,0,1}
\DefineNamedColor{named}{RoyalBlue}     {cmyk}{1,0.50,0,0}
\DefineNamedColor{named}{TealBlue}      {cmyk}{0.86,0,0.34,0.02}
\DefineNamedColor{named}{CarnationPink} {cmyk}{0,0.63,0,0}

\title{Scala, a Scalable Language} 
\author{Alan Dipert} 
\date{September 10th, 2009} 

\begin{document} 
\maketitle 

\begin{frame} 
\frametitle{What Scala Is}
\begin{columns}[c]
  \column{0.5in}
    \includegraphics[width=1.0in]{graphics/scala_logo.png} 
  \column{2.5in}
    \begin{itemize}
      \item<1-> Multi-paradigm: purely object oriented, but with functional capabilities
      \item<2-> Statically typed
      \item<3-> Seamless Java interoperability
      \item<4-> Open source (BSD License)
      \item<5-> Latest of several JVM languages from creator Martin Odersky and team at EPFL
      \item<6-> Cool: \tt\small{def fact(x:BigInt):BigInt = if (x==0) 1 else x * fact(x-1)}
    \end{itemize}
\end{columns}
\end{frame} 

\begin{frame} 
\frametitle{Object Oriented Programming in Scala}
\begin{columns}[c]
  \column{1.0in}
    \includegraphics[width=1.8in]{graphics/oop_dot.pdf} 
  \column{2.0in}
    \begin{itemize}
      \item<1-> Classes and objects similar to Java...
      \item<2-> ...except "everything is an object."  No simple types like Java
      \item<3-> Automatic getters/setters (like :attr\_accessor in Ruby)
      \item<4-> "Singleton objects" instead of static methods in classes
      \item<5-> "Traits," behaviors mixed in with classes and objects, simplify multiple inheritance
    \end{itemize}
\end{columns}
\end{frame} 

\begin{frame}[fragile]
\frametitle{Class Definition in Scala}
{\tt \tiny
\begin{semiverbatim}
\uncover<1->{//Person objects take the name 'n' and the age 'a' in the constructor}
\uncover<1->{class Person(n:String, a:Int) \{}
\uncover<2->{  //name and age are immutable instance values}
\uncover<2->{  val name = n}
\uncover<2->{  val age = a}
\uncover<3->{  //Person provides its own toString method that returns something sensible}
\uncover<3->{  override def toString = "My name is %s and I'm %d years old.".format(name,age)}
\uncover<1->{\}}
\uncover<4->{}
\uncover<4->{val thePrez = new Person("Barrack Obama", 48)}
\uncover<5->{}
\uncover<5->{println(thePrez toString)}
\uncover<5->{//prints "My name is Barrack Obama and I'm 48 years old."}
\end{semiverbatim}
}
\end{frame} 

\begin{frame}[fragile]
\frametitle{Inheritance in Scala}
{\tt \tiny
\begin{semiverbatim}
\uncover<1->{class Person(n:String, a:Int) \{}
\uncover<1->{  val name = n}
\uncover<1->{  val age = a}
\uncover<1->{  override def toString = "My name is %s and I'm %d years old.".format(name,age)}
\uncover<1->{\}}
\uncover<2->{}
\uncover<2->{class Student(u:String, n:String, a:Int) extends Person(n,a) \{}
\uncover<3->{  val school = u}
\uncover<3->{  override def toString = super.toString+" I go to %s.".format(school)}
\uncover<3->{\}}
\uncover<4->{}
\uncover<4->{val alan = new Student("RIT", "Alan Dipert", 25)}
\uncover<5->{}
\uncover<5->{println(alan toString)}
\uncover<5->{//prints "My name is Alan Dipert and I'm 25 years old. I go to RIT."}
\end{semiverbatim}
}
\end{frame} 

\begin{frame}[fragile]
\frametitle{Traits in Scala}
{\tt \tiny
\begin{semiverbatim}
\uncover<1->{class Person(n:String, a:Int) \{}
\uncover<1->{  val name = n}
\uncover<1->{  val age = a}
\uncover<1->{  override def toString = "My name is %s and I'm %d years old.".format(name,age)}
\uncover<1->{\}}
\uncover<1->{}
\uncover<1->{class Student(u:String, n:String, a:Int) extends Person(n,a) \{}
\uncover<1->{  val school = u}
\uncover<1->{  override def toString = super.toString+" I go to %s.".format(school)}
\uncover<1->{\}}
\uncover<2->{}
\uncover<2->{trait Cool \{}
\uncover<2->{  override def toString = super.toString+" And I'm cool."}
\uncover<2->{\}}
\uncover<2->{}
\uncover<3->{trait LaidBack \{}
\uncover<3->{  override def toString = super.toString+" And I'm laid back."}
\uncover<3->{\}}
\uncover<4->{}
\uncover<4->{val ferris = new Student("Shermer High", "Ferris Bueller", 17) with LaidBack with Cool}
\uncover<5->{}
\uncover<5->{println(ferris toString)}
\uncover<5->{//"My name is Ferris Bueller and I'm 17 years old. I go to Shermer High. And I'm laid back. And I'm cool."}
\end{semiverbatim}
}
\end{frame} 

\begin{frame} 
\frametitle{Object Oriented Programming Conclusions}
\begin{columns}[c]
  \column{1.0in}
    \includegraphics[width=1.5in]{graphics/ferris.jpg} 
  \column{2.0in}
    \begin{itemize}
      \item<1-> Scala's object system gives you a lot of control
      \item<2-> It feels much more like Ruby or Python than Java or C++
      \item<3-> It's a carefully designed fusion of useful features and Java compatibility
      \item<4-> Like a lot of Scala, you can start by doing things the Java way
    \end{itemize}
\end{columns}
\end{frame} 

\begin{frame} 
\frametitle{Functional Programming Introduction}
\begin{columns}[c]
  \column{0.5in}
    \includegraphics[width=1.0in]{graphics/hat.png} 
  \column{2.5in}
    \begin{itemize}
      \item<1-> FP is the paradigm of avoiding state and mutability in code, "functions and data the same"
      \item<2-> Around practically since Lisp (1958), popular lately because of focus on concurrency
      \item<3-> Side-effect free functions and immutable data structures have definite advantages in many domains
      \item<4-> On the JVM, FP alternatives are Clojure, Groovy
      \item<5-> In Java terms, anonymous inner classes on crack
    \end{itemize}
\end{columns}
\end{frame} 

\begin{frame} 
\frametitle{Functional Programming in Scala}
\begin{itemize}
  \item<1-> Scala has first-class functions
  \item<2-> Many Collection classes provide convenience higher order functions: \tt\small{reduceLeft, foldLeft, filter, forall}
  \item<3-> Lots of functional stand-bys: \tt\small{head, tail}
  \item<4-> Functional techniques help you avoid off-by one errors, temporary variables, and explicit loops
  \item<5-> Chained higher-order functions let you write common imperative for/if patterns concisely
\end{itemize}
\end{frame} 

\begin{frame} 
\frametitle{Functional Programming in Scala Examples}
\begin{itemize}
  \item<1-> Sum a list of Integers: \tt{List(1,5,3,2).reduceLeft((x,y) => x+y) (Int = 11)}
  \item<2-> More concisely: \tt{List(1,5,3,2).reduceLeft(\_+\_) (Int = 11)}
  \item<3-> Prune odd numbers: \tt{List(1,2,3,4).filter(x => if (x\%2 == 0) true else false) (List[Int] = List(2,4))}
  \item<4-> Prune odd numbers, get sum: \tt{List(1,2,3,4).filter(x => if (x\%2 == 0) true else false).reduceLeft(\_+\_) (Int = 6)}
\end{itemize}
\end{frame} 

\begin{frame} 
\frametitle{Putting it All Together: a Simple Application}
\begin{itemize}
  \item<1-> Three languages, one application. C, Java, and Scala
  \item<2-> Imperative, mostly imperative, and a mix of functional/imperative paradigms
  \item<3-> \textbf{Warning: always use the right tool for the job!}
  \item<4-> This is just a demonstration.
  \item<5-> Our program is supposed to take one argument, a comma-delimited list of integers to sum.
  \item<6-> Usage: \tt{./summer 1,2,3} should print 6
\end{itemize}
\end{frame} 

\begin{frame} 
\frametitle{Old Reliable: C}
\tt\tiny{
\input{code/summer.c.tex}
}
\end{frame} 

\begin{frame} 
\frametitle{The Clinton Years: Java}
\tt\tiny{
\input{code/Summer.java.tex}
}
\end{frame} 

\begin{frame} 
\frametitle{Today: Scala}
\tt\tiny{
\input{code/Summer.scala.tex}
}
\end{frame} 

\begin{frame} 
\frametitle{Thank You!}
\begin{itemize}
  \item<1-> \href{http://www.scala-lang.org/}{The Scala Homepage}
  \item<1-> \href{http://www.artima.com/shop/programming_in_scala}{Programming in Scala, the definitive book}
\end{itemize}
\begin{itemize}
  \item<1-> \href{mailto:alan@nuklei.com}{alan@nuklei.com}
  \item<1-> \href{http://alan.dipert.org/}{http://alan.dipert.org/}
\end{itemize}
\end{frame} 

\end{document} 
